Throughout, $C,C_1,C_2,\dots$ and $c_0$ denote constants that depend only on the sub-Gaussian constant $K_0$ and may change from line to line. We fix a class $c$ and drop it from the notation when no confusion arises.

\subsection{Proof of Proposition~\ref{prop:inv}}
The reference model $w$, the sketch $\Phi$ and the label $c$ are the same for all clients, so
$\mu_i^c=\int \Phi\nabla_w\ell(w;x,c)\,d\cP_i(x\mid Y=c)$ is determined by $\cP_i(X\mid Y=c)$ alone. Neither the proportion $\cP_i(Y=c)$ nor the number of samples enters. If $i$ and $j$ are compatible, then $\cP_i(X\mid Y=c)=\cP_j(X\mid Y=c)$ on each shared class, hence $\mu_i^c=\mu_j^c$ and $\rho^c_{ij}=1$. Per-sample clipping (Proposition~\ref{prop:dp}) is applied identically by all clients, so the same holds for the expected clipped signature. \qed

\subsection{Proof of Lemma~\ref{lem:pop}}
Write the half-means as $a_i=\mu_i+\zeta_{i1}$, $b_i=\mu_i+\zeta_{i2}$, where $\zeta_{i1},\zeta_{i2}$ are independent, zero-mean and have covariance $\Sigma_i/h_i$ (any covariance works, and it may differ across clients). Let $t_i=\tr\Sigma_i/h_i$. Then
$\E\langle a_i,b_i\rangle=a_i$ and $\E\|a_i\|^2=\E\|b_i\|^2=a_i+t_i$, so the population split-half reliability is $r_i^\star=a_i/(a_i+t_i)$ and
\[R_i^\star=\frac{2r_i^\star}{1+r_i^\star}=\frac{a_i}{a_i+t_i/2}.\]
For the full means, $\E\|m_i\|^2=a_i+t_i/2=a_i/R_i^\star$ and, by independence across clients, $\E\langle m_i,m_j\rangle=\langle\mu_i,\mu_j\rangle$. The population cosine is therefore
$\langle\mu_i,\mu_j\rangle\sqrt{R^\star_iR^\star_j}/\sqrt{a_ia_j}$. Dividing by $\sqrt{R_i^\star R_j^\star}$ gives $\rho_{ij}$ exactly. \qed

\subsection{Proof of Theorem~\ref{thm:conc}}
\emph{Noise representation.} A half-mean noise vector is $\zeta=(\Sigma/h)^{1/2}\bar z$ with $\bar z=h^{-1/2}\sum_{s}z_s$. Its coordinates are independent, zero-mean, unit-variance and $CK_0$-sub-Gaussian.

\emph{Statistics.} $\rhohat_{ij}$ and $R_i$ are functions of the nine statistics
$X_1=\langle a_i,b_i\rangle$, $X_2=\|a_i\|^2$, $X_3=\|b_i\|^2$, $X_4=\|m_i\|^2$ (and $X_5,\dots,X_8$ for $j$) and $X_9=\langle m_i,m_j\rangle$. We normalise them as $x_k=X_k/a_i$ for the statistics of $i$, $x_k=X_k/a_j$ for those of $j$, and $x_9=X_9/\sqrt{a_ia_j}$. Their expectations are $x^\star=(1,1+\kappa_i,1+\kappa_i,1+\kappa_i/2,\ \dots,\ \rho_{ij})$ with $\kappa_i=\tr\Sigma/(a_ih_i)\le\kappa$. Each statistic has the form
constant $+$ linear term $\langle\mu,\zeta\rangle$ $+$ quadratic or bilinear form in $\zeta$'s.

\emph{Linear terms.} $\langle\mu_i,\zeta\rangle$ is sub-Gaussian with variance proxy $\le CK_0^2\,\mu_i^\top\Sigma\mu_i/h\le CK_0^2a_i\|\Sigma\|_\op/h$. After normalisation by $a_i$ (or by $\sqrt{a_ia_j}$ for cross terms), with probability $1-\delta'$ it is at most
$C\sqrt{\|\Sigma\|_\op/(ah)}\sqrt{\log(2/\delta')}=C\sqrt{\kappa/r_s}\sqrt{\log(2/\delta')}$, because $\|\Sigma\|_\op=\sigma^2/r_s$.

\emph{Quadratic and bilinear terms.} $\|\zeta\|^2-\tr(\Sigma/h)$ and $\langle\zeta,\zeta'\rangle$ for independent $\zeta,\zeta'$ are centred quadratic forms $\bar Z^\top M\bar Z$ of the concatenated vector $\bar Z$, with $\|M\|_F\le\|\Sigma\|_F/h$ and $\|M\|_\op\le\|\Sigma\|_\op/h$. The Hanson--Wright inequality~\citep{rudelson2013hansonwright} gives, with probability $1-\delta'$, a deviation of at most $C K_0^2(\|M\|_F\sqrt{\log(2/\delta')}+\|M\|_\op\log(2/\delta'))$. After normalisation by $a$ this is at most
$C(\kappa/\sqrt{r_e})\sqrt{\log(2/\delta')}+C(\kappa/r_s)\log(2/\delta')$, using $\|\Sigma\|_F=\sigma^2/\sqrt{r_e}$. The second term is at most $C\eta^2\le Cc_0\eta$.

\emph{Union bound.} Taking $\delta'=\delta/6$ over the at most six independent noise terms that appear in each group of statistics, we obtain $\|x-x^\star\|_\infty\le C\eta$ with probability $1-\delta$.

\emph{Smoothness.} Write $\rhohat=F(x)$ with $r_i=x_1/\sqrt{x_2x_3}$, $R_i=2r_i/(1+r_i)$ (similarly for $j$), and $\rhohat=x_9/\sqrt{x_4x_8R_iR_j}$ before clipping. By Lemma~\ref{lem:pop}, $F(x^\star)=\rho_{ij}$. For $\kappa\le1$ and $\|x-x^\star\|_\infty\le Cc_0$ with $c_0$ small enough, the argument stays in a compact set on which $x_2,x_3,x_4\in[1/2,3]$, $r_i\ge1/5$ and $R_i\ge1/3$. On that set $F$ is continuously differentiable with gradient bounded by an absolute constant. Hence $|F(x)-F(x^\star)|\le C_1\eta$, and the same argument gives the bound for $R_i$. Clipping to $[-1,1]$ can only reduce the error, because $\rho_{ij}\in[-1,1]$. For the raw cosine, $x_9/\sqrt{x_4x_8}\to\rho_{ij}/\sqrt{(1/R^\star_i)(1/R^\star_j)}=\rho_{ij}\sqrt{R_i^\star R_j^\star}$. \qed

\subsection{Proof of Theorem~\ref{thm:rec}}
We use two facts, both on $\mathcal E$.

\emph{Claim 1 (conflicting clusters cannot be linked).} Let $A,B$ be conflict-free and suppose some $i\in A$, $j\in B$ conflict on class $c$. Since $A$ is conflict-free, all members of $A$ holding $c$ share one class-conditional $\cP_A^c$; likewise for $B$, with $\cP_A^c\neq\cP_B^c$. Hence \emph{every} pair in $A\times B$ holding $c$ conflicts on $c$. By Assumption~\ref{ass:sep} and $\mathcal E$, each such pair has $\rhohat\le1-\Delta+\gamma/2$. Thus $\mathrm{num}^c_{AB}\le(1-\Delta+\gamma/2)\,\mathrm{den}^c_{AB}$ and
\[\frac{\mathrm{num}^c_{AB}+\lambda}{\mathrm{den}^c_{AB}+\lambda}<\tau\iff \mathrm{den}^c_{AB}(\tau-1+\Delta-\gamma/2)>\lambda(1-\tau).\]
Since $\tau-1+\Delta\ge\gamma$, the left side is at least $\mathrm{den}^c_{AB}\gamma/2\ge W_{ij}^c\gamma/2>w_0\gamma/2=\lambda(1-\tau)$. The minimum over classes in~\eqref{eq:link} is therefore below $\tau$.

\emph{Claim 2 (compatible clusters are always linkable).} If $A\cup B$ is conflict-free, then every pair in $A\times B$ that holds a class $c$ has $\rho^c=1$ (Proposition~\ref{prop:inv}), so $\rhohat^c\ge1-\gamma/2$. Hence every class term of~\eqref{eq:link} is at least $1-\gamma/2\ge1-(1-\tau)/2=(1+\tau)/2>\tau$; classes without evidence contribute $1$. So $L(A,B)>\tau$ whenever the shared evidence is at least $e_{\min}$.

(a) Singletons are conflict-free. By Claim~1 a merge is only performed between clusters whose union is conflict-free, so the invariant is preserved by every merge, in any order. (b) At termination every pair of clusters with shared evidence $\ge e_{\min}$ has $L<\tau$; by Claim~2 their union is not conflict-free. (c) Under full support, two clients of different groups $k\ne l$ both hold a class on which $\cP^{(k)}$ and $\cP^{(l)}$ differ, so they conflict. By (a) each output cluster is therefore contained in a single group. Two output clusters inside the same group have a conflict-free union and share evidence $\ge e_{\min}$, which contradicts (b). Each group thus forms exactly one cluster. \qed

\subsection{Proof of Corollary~\ref{cor:sc}}
There are at most $N^2C/2$ pairwise class estimates and $NC$ reliabilities; we apply Theorem~\ref{thm:conc} to each with $\delta/(N^2C)$ and write $\mathsf L=\log(12N^2C/\delta)$. The requirement $C_1\eta\le\gamma/2$ holds if $\sqrt{\kappa/r_s}\sqrt{\mathsf L}\le\gamma/(4C_1)$ and $(\kappa/\sqrt{r_e})\sqrt{\mathsf L}\le\gamma/(4C_1)$. With $\kappa=\sigma^2/(ah)$, these are the two terms of~\eqref{eq:hstar}; the side conditions $\kappa\le1$ and $\eta\le c_0$ are absorbed into $C_2$ and the additional requirement $h\ge\sigma^2/a$. For the evidence condition: on the event, $R\ge R^\star-\gamma/2\ge2/3-\gamma/2$ (since $\kappa\le1$). For $\gamma\le1/3$ this gives $W\ge1/4$, so $W>w_0$ holds whenever $\lambda<\gamma/(8(1-\tau))$. None of these conditions involves $|G_k|$. \qed

\subsection{Proof of Theorem~\ref{thm:lb}}
Let $e_1$ be a top eigenvector of $\Sigma$ ($\Sigma e_1=\|\Sigma\|_\op e_1$) and $e_2\perp e_1$ a unit vector. With $s^2=\Delta/2$ put
$u=\sqrt a(\sqrt{1-s^2}\,e_2+s\,e_1)$ and $v=\sqrt a(\sqrt{1-s^2}\,e_2-s\,e_1)$. Then $\|u\|^2=\|v\|^2=a$, $\cos(u,v)=1-2s^2=1-\Delta$ and $u-v=2\sqrt a\,s\,e_1$. Under $H_0$ the $n$ samples of client $j$ are i.i.d.\ $\mathcal N(u,\Sigma)$; under $H_1$ they are $\mathcal N(v,\Sigma)$. Moreover $\mu_i=u$ is known. Then
\[\mathrm{KL}\big(\mathcal N(u,\Sigma)^{\otimes n}\,\|\,\mathcal N(v,\Sigma)^{\otimes n}\big)=\frac n2(u-v)^\top\Sigma^{-1}(u-v)=\frac n2\cdot\frac{4as^2}{\|\Sigma\|_\op}=\frac{n\,a\,\Delta\,r_s}{\sigma^2}.\]
By the Bretagnolle--Huber inequality (\citealp[Thm.~2.2]{tsybakov2009}), every test $\psi$ satisfies $\max_{H\in\{0,1\}}\Pr_H(\psi\neq H)\ge\frac14\exp(-\mathrm{KL})$. A procedure that sees only half-means (as DisCo-CFL does), or that does not know $\mu_i$, has access to less information, so the bound applies to it as well. \qed

\subsection{Proof of Proposition~\ref{prop:dp}}
(i) Replacing one record by another of the same label changes one clipped summand, hence the half-mean by at most $2C_{\rm clip}/h$ in $\ell_2$. The added noise has standard deviation $\sigma_{\rm dp}$ times this sensitivity, so the analytic Gaussian mechanism~\citep{balle2018analytic} gives $(\varepsilon,\delta)$-DP for that vector. Every record enters exactly one released vector, so by parallel composition the whole signature is $(\varepsilon,\delta)$-DP. The server's computations are post-processing~\citep{dwork2014dp}. The class supports and half sizes are released and are not protected. (ii) The DP noise is zero-mean and independent of everything else. The proof of Lemma~\ref{lem:pop} used only these two properties, so it still holds with $t_i=\tr\Sigma_{\rm clip}/h+p(2\sigma_{\rm dp}C_{\rm clip}/h)^2$. In the proof of Theorem~\ref{thm:conc}, the Gaussian DP component adds a linear term with normalised standard deviation $2\sigma_{\rm dp}C_{\rm clip}/(h\sqrt a)$ and a quadratic term with normalised Frobenius scale $4\sqrt p\,\sigma^2_{\rm dp}C^2_{\rm clip}/(h^2a)$. Requiring each to be at most $\gamma/(4C_1\sqrt{\mathsf L})$ gives the two stated conditions. (iii) This follows from the last statement of Theorem~\ref{thm:conc} with the inflated $t_i$. \qed

\subsection{Proof of Proposition~\ref{prop:expl}}
Output clusters are conflict-free (Theorem~\ref{thm:rec}(a)), so the members of $A$ holding $c$ share $\cP_A^c$, and likewise for $B$. If $\cP_A^c\neq\cP_B^c$, the computation of Claim~1 with $\mathrm{den}^c_{AB}>w_0$ gives a pooled value below $\tau$, so $c\in\hat D_{AB}$. Otherwise all pairs have $\rho^c=1$ and the pooled value is at least $(1+\tau)/2>\tau$ (Claim~2), so $c\notin\hat D_{AB}$. \qed
